\providecommand{\FloatBarrier}{\clearpage}

\section{Performance of the Bayesian Classifier}
\label{sec:performance}

\subsection{Validation Metrics}
\label{sec:validation-metrics}

We evaluate the classifiers using the matched COSMOS2020/FARMER reference
sample \citep{Weaver2022}.  The one-to-one match, constructed with a maximum
separation of $0.5\arcsec$, contains 193,586 objects: 8,788 reference-labeled
stars and 184,798 reference-labeled galaxies.  Throughout this section, $r$
denotes the Rubin $r$-band CModel magnitude.

The matched sample contains about 21 times more galaxies than stars, and the
imbalance is stronger at faint magnitudes.  A single accuracy value would
therefore be dominated by the much larger galaxy population.  We instead
evaluate the star-candidate and galaxy-candidate samples separately.  For a
target class $X$, where $X$ is either stars or galaxies, we define
completeness, purity, and contamination as

\begin{align}
    C_X &\equiv \frac{\mathrm{TP}_X}
                       {\mathrm{TP}_X+\mathrm{FN}_X},
    \label{eq:completeness}\\
    P_X &\equiv \frac{\mathrm{TP}_X}
                       {\mathrm{TP}_X+\mathrm{FP}_X},
    \label{eq:purity}\\
    f_{\mathrm{cont},X} &\equiv 1-P_X
      =\frac{\mathrm{FP}_X}{\mathrm{TP}_X+\mathrm{FP}_X}.
    \label{eq:contamination}
\end{align}

Completeness is the fraction of reference objects of class $X$ that are
recovered, whereas purity is the fraction of the selected sample that belongs
to class $X$ according to the reference labels.  Purity is particularly
important for the imbalanced faint sample because even a modest false-positive
rate can produce substantial contamination \citep{SaitoRehmsmeier2015}.

For a receiver operating characteristic (ROC) curve, the false-positive rate
for target class $X$ is

\begin{equation}
    \mathrm{FPR}_X \equiv
    \frac{\mathrm{FP}_X}{\mathrm{FP}_X+\mathrm{TN}_X}.
    \label{eq:false-positive-rate}
\end{equation}

Its denominator contains all reference members of the opposite class and
therefore differs from the denominator of contamination.  The true-positive
rate on the ROC curve is equal to $C_X$.  The area under the ROC curve (AUC)
summarizes how well a continuous score orders the two classes
\citep{Fawcett2006ROC}.

\subsection{Single-band Morphology Performance}
\label{sec:single-band-performance}

We use the continuous $r$-band scores to compare how well the two reference
classes are separated, and we use the fixed catalog rules to measure the
completeness and purity of the selected samples.  Every probabilistic catalog
classification in this section uses $p_{{\rm S},r}\geq0.5$ for a star-like
assignment.  The Rubin comparison uses its default boundary,
$\Delta m_r\leq0.016\,\mathrm{mag}$.

\begin{figure*}[!t]
    \centering
    \includegraphics[width=\textwidth]{fig09_morphology_only_roc.png}
    \caption{Performance of the continuous $r$-band morphology scores for the
    $0.5\arcsec$ COSMOS validation sample.  The left and right panels treat stars
    and galaxies as the target class, respectively.  The gray dash-dotted curve
    uses raw $\Delta m_r$, the blue curve uses the fitted morphology distributions
    without the source-count prior, and the green curve includes the prior.  The
    red diamond marks the fixed Rubin \texttt{extendedness} assignment at
    $\Delta m_r=0.016\,\mathrm{mag}$, and the black circle marks the fixed
    probabilistic assignment at $p_{{\rm S},r}=0.5$.  Values in parentheses give
    the AUC of each continuous score.}
    \label{fig:rband-roc}
\end{figure*}

\subsubsection{Ranking of the Continuous Scores}
\label{sec:rband-ranking}

Figure~\ref{fig:rband-roc} compares three continuous $r$-band morphology
scores.  The first is the measured quantity

\begin{equation}
    \Delta m_r=m_{{\rm PSF},r}-m_{{\rm CModel},r},
    \label{eq:delta-m-r}
\end{equation}

which underlies the Rubin \texttt{extendedness} classification.  The gray
curve therefore measures the discriminatory information in the raw
$\Delta m_r$ values.  The fixed Rubin catalog classification is shown
separately by the red diamond at $\Delta m_r=0.016\,\mathrm{mag}$.  Likewise,
the fixed probabilistic classification at $p_{{\rm S},r}=0.5$ is marked on the
prior-informed score by the black circle.

The blue and green curves use the fitted class-conditional morphology
distributions described in Section~\ref{sec:mixture-model}.  The blue score
omits the magnitude-dependent source-count prior, whereas the green score
includes the prior from Section~\ref{sec:sg-priors}.  All three scores are
evaluated on the same 193,516 objects with finite $r$-band morphology
measurements.

The AUC is 0.871 for raw $\Delta m_r$, 0.899 for the morphology model without
the count prior, and 0.901 after the prior is included.  The modeled
class-conditional distributions therefore improve the overall ordering, while
the count prior changes the global AUC only slightly.  The ROC curves cross:
raw $\Delta m_r$ lies above the modeled scores over parts of the displayed
false-positive-rate range, while the modeled scores perform better elsewhere.
The higher AUC therefore describes an improvement averaged over the full
curve, not uniform dominance at every operating region.

\subsubsection{Completeness and Purity of the Fixed Catalog Samples}
\label{sec:rband-catalog-performance}

Figure~\ref{fig:rband-fixed-performance} compares the samples produced by the
fixed $p_{{\rm S},r}=0.5$ and Rubin \texttt{extendedness} rules in three
magnitude intervals.  The unresolved-selected and resolved-selected samples
are treated as star-candidate and galaxy-candidate samples, respectively.  We
show completeness and purity; contamination is simply $1-P_X$.

At $16\leq r<24$, the probabilistic star-candidate sample has completeness
0.878 and purity 0.975, similar to the Rubin selection in this interval.  At
$24\leq r<25$, the two methods make a different trade-off: the probabilistic
selection increases stellar purity from 0.318 to 0.558, while stellar
completeness decreases from 0.582 to 0.268.

At $25\leq r<26$, no evaluated object satisfies $p_{{\rm S},r}\geq0.5$.
Stellar completeness is therefore zero, and stellar purity is undefined
because the selected sample is empty.  Thus, the fixed prior-informed
$r$-band rule does not produce a usable star-candidate sample in the faintest
interval.  The galaxy-candidate sample remains nearly complete and highly
pure, but these values are aided by the overwhelming numerical dominance of
galaxies and do not imply comparable recovery of the rare stellar population.

\begin{figure*}[!t]
    \centering
    \includegraphics[width=\textwidth]{fig10_single_band_performance.png}
    \caption{Completeness (top row) and purity (bottom row) of the
    star-candidate sample (left column) and galaxy-candidate sample (right column)
    in three $r$-band magnitude intervals.  Horizontally offset markers compare
    the fixed $p_{{\rm S},r}=0.5$ assignment with the default Rubin
    \texttt{extendedness} assignment; the magnitude intervals are not connected
    by lines.  Error bars show 68\% Wilson binomial intervals, and annotations give
    the number of selected objects.  At $25\leq r<26$, the probabilistic method
    selects no star-like objects, so stellar completeness is zero and stellar
    purity is undefined and is not plotted.}
    \label{fig:rband-fixed-performance}
\end{figure*}

\subsection{Multiband Morphology and Color Information}
\label{sec:multiband-color}

Morphology and broadband colors provide complementary information for
star--galaxy separation
\citep{Fadely2012MultiBandSG,Kim2015HybridSG,SevillaNoarbe2018DESY1}.
Figure~\ref{fig:multiband-color-roc} compares scores based on $r$-band
morphology, $gri$ morphology, all six $ugrizy$ morphology measurements, and a
combined morphology-plus-color score.  The band-level morphology information
is combined before the magnitude-dependent source-count prior is applied.
Within each magnitude interval, all curves use the same complete-case sample.

\begin{figure*}[!t]
    \centering
    \includegraphics[width=\textwidth]{fig11_multiband_color_roc.png}
    \caption{Ranking performance obtained from multiband morphology and color
    information.  The top and bottom rows treat stars and galaxies as the target
    class, respectively, and the columns show three $r$-band magnitude intervals.
    Curves compare morphology information from the $r$, $gri$, and $ugrizy$ bands
    with the combined $ugrizy$ morphology-plus-color score.  Every panel uses a
    common complete-case sample.  The red diamond marks the default Rubin
    $r$-band \texttt{extendedness} assignment at
    $\Delta m_r=0.016\,\mathrm{mag}$.  Values in parentheses give the AUC.}
    \label{fig:multiband-color-roc}
\end{figure*}

The combined score supplements the $ugrizy$ morphology information with the
output of a Random Forest classifier based on the five adjacent colors
$(u-g,g-r,r-i,i-z,z-y)$ \citep{Breiman2001RandomForests}.  Figure~\ref{fig:multiband-color-roc}
uses this output together with the multiband morphology score as a continuous
classifier score.

% Before submission, verify from the implementation that the Random Forest
% values used here are held-out or out-of-fold predictions and confirm the
% exact rule used to combine the color and morphology scores.

For the star-positive comparison, the $ugrizy$ morphology AUC is 0.978,
0.777, and 0.721 in the three magnitude intervals, compared with 0.965,
0.759, and 0.734 for $r$-band morphology alone.  The six-band morphology score
therefore gives modest gains in the first two intervals, while the $r$-only
score is slightly higher in the faintest interval.  The $gri$ score is also
slightly below the $r$-only result in the two faint intervals.  These
comparisons show that the contribution of an additional morphology band
depends on the useful depth of its measurements, rather than simply on the
number of bands included.

Including the color score changes the $ugrizy$ morphology AUC from 0.978 to
0.986, from 0.777 to 0.815, and from 0.721 to 0.779.  The corresponding
increments are $+0.008$, $+0.038$, and $+0.058$, so the additional separation
provided by color becomes progressively larger across the three displayed
magnitude intervals.  The top row of
Figure~\ref{fig:multiband-color-roc} gives the completeness of reference stars
as galaxies are admitted into the star-selected sample, while the bottom row
gives the completeness of reference galaxies as stars are admitted into the
galaxy-selected sample.  The two rows therefore expose the trade-off relevant
to each selected class directly.

\subsection{Color Distributions of the Classifier-selected Samples}
\label{sec:classifier-selected-colors}

Figure~\ref{fig:color-distributions} shows how the unresolved and resolved
assignments from the combined $ugrizy$ morphology-plus-color classifier are
distributed in four two-dimensional color projections.  This provides a
direct view of the color-space partition produced by the final classifier.

\begin{figure*}[!t]
    \centering
    \includegraphics[width=\textwidth,height=0.64\textheight,keepaspectratio]{fig12_ugrizy_color_method_labels.png}
    \caption{Color--color distributions of the samples assigned to the unresolved
    and resolved classes by the combined $ugrizy$ morphology-plus-color
    classifier, shown for $16\leq r<25$ (left) and $25\leq r<26$ (right).  Blue
    points show unresolved assignments, and red contours show the resolved
    population.  The two intervals contain 5,855 and 180 unresolved assignments
    and 106,740 and 25,707 resolved assignments, respectively.  The overlap in
    the individual two-dimensional projections shows that the final assignment
    cannot be represented by a single color--color boundary.}
    \label{fig:color-distributions}
\end{figure*}

Classifier-selected unresolved objects are plotted as blue points, while red
density contours show the classifier-selected resolved population.  The
$16\leq r<25$ interval contains 5,855 unresolved and 106,740 resolved
assignments; the $25\leq r<26$ interval contains 180 and 25,707, respectively.
The small number of faint unresolved assignments accounts for the sparse blue
sample in the right column.  In each individual projection, the blue points
and red contours overlap substantially.  The final classification is
therefore not equivalent to a simple boundary in any one of these
color--color planes; it uses their joint information together with
morphology.

Taken together, the results show that modeling the $r$-band morphology
distributions improves the global AUC relative to raw $\Delta m_r$, although
the ROC curves cross and the gain is not uniform over all false-positive
rates.  At the fixed $p_{{\rm S},r}=0.5$ boundary, the prior-informed $r$-band
selection becomes increasingly conservative and selects no star-like objects
in the faintest interval.  In the multiband comparison, the incremental AUC
gain from color increases toward fainter magnitudes within the evaluated
sample.

\FloatBarrier